\chapter{2006年辽宁卷（理）}
\section{单选题}
\begin{problem}
设集合 \(A = \{1, 2\}\)，则满足 \(A \cup B = \{1, 2, 3\}\) 的集合 \(B\) 的个数是（\quad）

\choices
    {\(1\)}
    {\(3\)}
    {\(4\)}
    {\(8\)}
\end{problem}

\begin{answer}
C
\end{answer}

\begin{solution}
因为 \(A\cup B=\{1,2,3\}\)，所以 \(B\) 必须含有 \(3\)，且不能含有 \(1\)，\(2\)，\(3\) 以外的元素. 集合 \(B\) 是否含有 \(1\)，\(2\) 分别有两种选择，因此共有 \(2^2=4\) 个集合，故选 C.
\end{solution}

\begin{problem}
设 \( f(x) \) 是 \( \R \) 上的任意函数，则下列叙述正确的是（\quad）

\choices
    {\(f(x)f(-x)\) 是奇函数}
    {\(f(x)|f(-x)|\) 是奇函数}
    {\(f(x) - f(-x)\) 是偶函数}
    {\(f(x) + f(-x)\) 是偶函数}
\end{problem}

\begin{answer}
D
\end{answer}

\begin{solution}
令 \(h(x)=f(x)+f(-x)\)，则 \(h(-x)=f(-x)+f(x)=h(x)\)，所以 \(h(x)\) 是偶函数. 对于 \(f(x)f(-x)\)，有 \(f(-x)f(x)=f(x)f(-x)\)，因此它总是偶函数，但取 \(f(x)=1\) 时不是奇函数. 对于 \(f(x)-f(-x)\)，有 \(f(-x)-f(x)=-\bigl(f(x)-f(-x)\bigr)\)，因此它总是奇函数，但取 \(f(x)=x\) 时不是偶函数. 最后，取 \(f(x)=1\) 时 \(f(x)|f(-x)|=1\) 为偶函数，取 \(f(x)=x\) 时 \(f(x)|f(-x)|=x|x|\) 为奇函数，所以该函数的奇偶性不能由任意的 \(f\) 确定. 故选 D.
\end{solution}

\begin{problem}
给出下列四个命题：

\begin{circlelist}
\item 垂直于同一直线的两条直线互相平行；
\item 垂直于同一平面的两个平面互相平行；
\item 若直线 \(l_{1}\)，\(l_{2}\) 与同一平面所成的角相等，则 \(l_{1}\)，\(l_{2}\) 互相平行；
\item 若直线 \(l_{1}\)，\(l_{2}\) 是异面直线，则与 \(l_{1}\)，\(l_{2}\) 都相交的两条直线是异面直线.
\end{circlelist}

其中假命题的个数是（\quad）

\choices
    {\(1\)}
    {\(2\)}
    {\(3\)}
    {\(4\)}
\end{problem}

\begin{answer}
D
\end{answer}

\begin{solution}
第一个命题不成立. 例如取 \(l\colon x=y=0\)，\(l_{1}\colon y=z=0\) 与 \(l_{2}\colon x=0\)，\(z=1\)，\(l_{1}\) 和 \(l_{2}\) 都与 \(l\) 相交且垂直，但 \(l_{1}\) 与 \(l_{2}\) 是异面直线. 第二个命题不成立：平面 \(x=0\) 和 \(y=0\) 都垂直于平面 \(z=0\)，但它们沿 \(z\) 轴相交. 第三个命题不成立：\(z=0\) 平面内的 \(x\) 轴与 \(y\) 轴相交，且二者与该平面所成的角都为 \(0\)，但不平行. 第四个命题也不成立：取异面直线 \(l_{1}\)，\(l_{2}\)，在 \(l_{1}\) 上取一点 \(P\)，在 \(l_{2}\) 上取两个不同的点 \(Q_{1}\)，\(Q_{2}\)，则直线 \(PQ_{1}\)，\(PQ_{2}\) 分别与 \(l_{1}\)，\(l_{2}\) 相交，并且它们在 \(P\) 处相交，所以不是异面直线. 因此四个命题均为假命题，假命题的个数为 \(4\)，故选 D.
\end{solution}

\begin{problem}
双曲线 \( x^{2}-y^{2}=4 \) 的两条渐近线与直线 \(x=3\) 围成一个三角形区域，表示该区域的不等式组是（\quad）

\choices
    {\(\begin{cases}
x - y \geqslant 0,\\
x + y \geqslant 0,\\
0 \leqslant x \leqslant 3
\end{cases}\)}
    {\(\begin{cases}
x - y \geqslant 0,\\
x + y \leqslant 0,\\
0 \leqslant x \leqslant 3
\end{cases}\)}
    {\(\begin{cases}
x - y \leqslant 0,\\
x + y \leqslant 0,\\
0 \leqslant x \leqslant 3
\end{cases}\)}
    {\(\begin{cases}
x - y \leqslant 0,\\
x + y \geqslant 3,\\
0 \leqslant x \leqslant 3
\end{cases}\)}
\end{problem}

\begin{answer}
A
\end{answer}

\begin{solution}
由 \(x^2-y^2=4\) 得双曲线的两条渐近线为 \(y=x\) 与 \(y=-x\). 这两条直线和 \(x=3\) 围成的三角形顶点为 \((0,0)\)，\((3,3)\)，\((3,-3)\)，其内部满足 \(y\leqslant x\)，\(y\geqslant -x\)，\(0\leqslant x\leqslant3\)，即 \(x-y\geqslant0\)，\(x+y\geqslant0\)，\(0\leqslant x\leqslant3\). 故选 A.
\end{solution}

\begin{problem}
设 \(\oplus\) 是 \(\R\) 上的一个运算，\(A\) 是 \(\R\) 的非空子集，若对任意 \(a\)，\(b\in A\)，有 \(a\oplus b\in A\)，则称 \(A\) 对运算 \(\oplus\) 封闭. 下列数集对加法，减法，乘法和除法（除数不等于零）四则运算都封闭的是（\quad）

\choices
    {自然数集}
    {整数集}
    {有理数集}
    {无理数集}
\end{problem}

\begin{answer}
C
\end{answer}

\begin{solution}
自然数集对减法不封闭，整数集对除法不封闭，无理数集对加法或乘法也不封闭. 例如无理数 \(\sqrt{2}\) 与 \(-\sqrt{2}\) 的和为有理数 \(0\). 有理数集对加法、减法、乘法都封闭，且两个有理数相除（除数不为零）仍为有理数，故选 C.
\end{solution}

\begin{problem}
\(\triangle ABC\) 的三内角 \(A\)，\(B\)，\(C\)，所对边的长分别为 \(a\)，\(b\)，\(c\)，设向量 \(\bs{p} = (a + c, b)\)，\(\bs{q} = (b - a, c - a)\). 若 \(\bs{p} \parallel \bs{q}\)，则角 \(C\) 的大小为（\quad）

\choices
    {\(\frac{\pi}{6}\)}
    {\(\frac{\pi}{3}\)}
    {\(\frac{\pi}{2}\)}
    {\(\frac{2\pi}{3}\)}
\end{problem}

\begin{answer}
B
\end{answer}

\begin{solution}
由两向量平行的坐标条件，得 \((a+c)(c-a)-b(b-a)=0\)，即 \(c^2=a^2+b^2-ab\). 由余弦定理，\(c^2=a^2+b^2-2ab\cos C\)，所以 \(2ab\cos C=ab\). 由于三角形边长 \(a\)，\(b\) 均为正数，故 \(\cos C=\frac12\)，从而 \(C=\frac{\pi}{3}\)，故选 B.
\end{solution}

\begin{problem}
与方程 \( y = \e^{2x} - 2\e^{x} + 1 \)（\(x \geqslant 0\)）的曲线关于直线 \( y = x \) 对称的曲线方程（\quad）

\choices
    {\( y = \ln (1 + \sqrt{x}) \)}
    {\( y = \ln (1 - \sqrt{x}) \)}
    {\( y = -\ln (1 + \sqrt{x}) \)}
    {\( y = -\ln (1 - \sqrt{x}) \)}
\end{problem}

\begin{answer}
A
\end{answer}

\begin{solution}
原曲线方程可写为 \(y=(\e^x-1)^2\)，且 \(x\geqslant0\)，所以 \(\e^x-1\geqslant0\). 关于 \(y=x\) 对称后交换 \(x\)，\(y\)，得 \(x=(\e^y-1)^2\)，于是 \(\sqrt{x}=\e^y-1\)，从而 \(y=\ln(1+\sqrt{x})\)，故选 A.
\end{solution}

\begin{problem}
曲线 \(\frac{x^2}{10 - m} + \frac{y^2}{6 - m} = 1 \)（\(m < 6\)）与曲线 \(\frac{x^2}{5 - n} + \frac{y^2}{9 - n} = 1 \)（\(5 < n < 9\)）的（\quad）

\choices
    {焦距相等}
    {离心率相等}
    {焦点相同}
    {准线相同}
\end{problem}

\begin{answer}
A
\end{answer}

\begin{solution}
因为 \(m<6\)，第一条曲线是焦点在 \(x\) 轴上的椭圆，半长轴平方与半短轴平方之差为 \((10-m)-(6-m)=4\)，故其焦半距为 \(2\)，焦距为 \(4\). 因为 \(5<n<9\)，第二条曲线可写成 \(\frac{y^2}{9-n}-\frac{x^2}{n-5}=1\)，是焦点在 \(y\) 轴上的双曲线，焦半距的平方为 \((9-n)+(n-5)=4\)，焦距也为 \(4\). 故选 A.
\end{solution}

\begin{problem}
在等比数列 \(\{a_{n}\}\) 中，\(a_{1} = 2\)，前 \(n\) 项和为 \(S_{n}\)，若数列 \(\{a_{n} + 1\}\) 也是等比数列，则 \(S_{n}\) 等于（\quad）

\choices
    {\(2^{n + 1} - 2\)}
    {\(3n\)}
    {\(2n\)}
    {\(3^n -1\)}
\end{problem}

\begin{answer}
C
\end{answer}

\begin{solution}
设等比数列 \(\{a_n\}\) 的公比为 \(q\). 由 \(a_1=2\)，有 \(a_2=2q\)，\(a_3=2q^2\). 因为 \(\{a_n+1\}\) 也是等比数列，所以 \((2q+1)^2=3(2q^2+1)\)，即 \(2(q-1)^2=0\)，从而 \(q=1\). 因此 \(a_n=2\)，\(S_n=2n\)，故选 C.
\end{solution}

\begin{problem}
直线 \( y = 2k \) 与曲线 \( 9k^2 x^2 + y^2 = 18k^2 |x| \)（\(k \in \R\)，且 \(k \neq 0\)）的公共点的个数为（\quad）

\choices
    {\(1\)}
    {\(2\)}
    {\(3\)}
    {\(4\)}
\end{problem}

\begin{answer}
D
\end{answer}

\begin{solution}
将 \(y=2k\) 代入曲线方程，并利用 \(k\neq0\)，得 \(9x^2-18|x|+4=0\). 令 \(t=|x|\geqslant0\)，则 \(9t^2-18t+4=0\)，解得 \(t=1\pm\frac{\sqrt{5}}{3}\). 这两个值均为正数，各自对应 \(x=\pm t\) 两个不同的值，且 \(y=2k\) 固定，所以公共点共有 \(4\) 个，故选 D.
\end{solution}

\begin{problem}
已知函数 \( f(x) = \frac{1}{2} (\sin x + \cos x) - \frac{1}{2} |\sin x - \cos x| \)，则 \( f(x) \) 的值域是（\quad）

\choices
    {\( \left[-1,1\right] \)}
    {\( \left[-\frac{\sqrt{2}}{2},1\right] \)}
    {\( \left[-1,\frac{\sqrt{2}}{2}\right] \)}
    {\(\left[-1, -\frac{\sqrt{2}}{2}\right]\)}
\end{problem}

\begin{answer}
C
\end{answer}

\begin{solution}
当 \(\sin x\geqslant\cos x\) 时，\(|\sin x-\cos x|=\sin x-\cos x\)，从而 \(f(x)=\cos x\)；当 \(\sin x<\cos x\) 时，\(|\sin x-\cos x|=\cos x-\sin x\)，从而 \(f(x)=\sin x\)，所以 \(f(x)=\min\{\sin x,\cos x\}\). 由 \(-1\leqslant\sin x\leqslant1\) 与 \(-1\leqslant\cos x\leqslant1\)，得 \(f(x)\geqslant-1\)，且在 \(x=\frac{3\pi}{2}\) 时取到 \(-1\). 另外，\(f(x)\leqslant\frac{\sin x+\cos x}{2}=\frac{\sqrt{2}}{2}\sin\left(x+\frac{\pi}{4}\right)\leqslant\frac{\sqrt{2}}{2}\)，并且在 \(x=\frac{\pi}{4}\) 时取到 \(\frac{\sqrt{2}}{2}\). 函数 \(f\) 连续，根据介值定理，它取遍两个极值之间的所有值，故值域为 \(\left[-1,\frac{\sqrt{2}}{2}\right]\)，选 C.
\end{solution}

\begin{problem}
设 \(O(0,0)\)，\(A(1,0)\)，\(B(0,1)\)，点 \(P\) 是线段 \(AB\) 上的一个动点，\(\overrightarrow{AP} = \lambda \overrightarrow{AB}\)，若 \(\overrightarrow{OP} \cdot \overrightarrow{AB} \geqslant \overrightarrow{PA} \cdot \overrightarrow{PB}\)，则实数 \(\lambda\) 的取值范围是（\quad）

\choices
    {\(\frac{1}{2} \leqslant \lambda \leqslant 1\)}
    {\(1 - \frac{\sqrt{2}}{2} \leqslant \lambda \leqslant 1\)}
    {\(\frac{1}{2} \leqslant \lambda \leqslant 1 + \frac{\sqrt{2}}{2}\)}
    {\(1 - \frac{\sqrt{2}}{2} \leqslant \lambda \leqslant 1 + \frac{\sqrt{2}}{2}\)}
\end{problem}

\begin{answer}
B
\end{answer}

\begin{solution}
因为 \(P\) 在线段 \(AB\) 上，设 \(P=(1-\lambda,\lambda)\)，则 \(0\leqslant\lambda\leqslant1\)，且 \(\overrightarrow{AB}=(-1,1)\)，\(\overrightarrow{PA}=(\lambda,-\lambda)\)，\(\overrightarrow{PB}=(\lambda-1,1-\lambda)\). 于是 \(\overrightarrow{OP}\cdot\overrightarrow{AB}=2\lambda-1\)，\(\overrightarrow{PA}\cdot\overrightarrow{PB}=-2\lambda(1-\lambda)\). 题设不等式化为 \(2\lambda^2-4\lambda+1\leqslant0\)，即 \(1-\frac{\sqrt{2}}{2}\leqslant\lambda\leqslant1+\frac{\sqrt{2}}{2}\). 与 \(0\leqslant\lambda\leqslant1\) 取交集，得 \(1-\frac{\sqrt{2}}{2}\leqslant\lambda\leqslant1\)，故选 B.
\end{solution}

\section{填空题}
\begin{problem}
设 \( g(x) = \begin{cases}
\e^x, & x \leqslant 0,\\
\ln x, & x > 0,
\end{cases} \) 则 \( g\left(g\left(\frac{1}{2}\right)\right) = \)\fillinblank{}.
\end{problem}

\begin{answer}
\(\frac{1}{2}\)
\end{answer}

\begin{solution}
因为 \(\frac12>0\)，所以 \(g\left(\frac12\right)=\ln\frac12=-\ln2<0\). 因此 \(g\left(g\left(\frac12\right)\right)=g(-\ln2)=\e^{-\ln2}=\frac12\).
\end{solution}

\begin{problem}
\(\displaystyle\lim_{n\to\infty}\frac{\left(\frac{4}{5} - \frac{6}{7}\right) + \left(\frac{4}{5^2} - \frac{6}{7^2}\right) + \cdots + \left(\frac{4}{5^n} - \frac{6}{7^n}\right)}{\left(\frac{4}{6} - \frac{4}{5}\right) + \left(\frac{4}{6^2} - \frac{4}{5^2}\right) + \cdots + \left(\frac{4}{6^n} - \frac{4}{5^n}\right)} =\)\fillinblank{}.
\end{problem}

\begin{answer}
0
\end{answer}

\begin{solution}
记分子为 \(N_n\)，由等比数列求和公式，\(N_n=(1-5^{-n})-(1-7^{-n})=7^{-n}-5^{-n}\)，所以 \(N_n\to0\). 记分母为 \(D_n\)，则 \(D_n=4\left[\frac15(1-6^{-n})-\frac14(1-5^{-n})\right]\)，从而 \(D_n\to-\frac15\neq0\). 因此原极限为 \(\displaystyle\lim_{n\to\infty}\frac{N_n}{D_n}=0\).
\end{solution}

\begin{problem}
\(5\)名乒乓球队队员中，有\(2\)名老队员和\(3\)名新队员. 从中选出\(3\)名队员排成\(1\)，\(2\)，\(3\)号参加团体比赛，则入选的\(3\)名队员中至少有\(1\)名老队员，且\(1\)，\(2\)号中至少有\(1\)名新队员的排列方法有\fillinblank{}种.（以数作答）
\end{problem}

\begin{answer}
\(48\)
\end{answer}

\begin{solution}
按入选队员中老队员的人数分类. 当有 \(1\) 名老队员和 \(2\) 名新队员时，先选老队员有 \(2\) 种，选出的 \(2\) 名新队员排入另外两个位置有 \(\mathrm A_3^2\) 种，老队员有 \(3\) 个位置可放，共有 \(2\cdot\mathrm A_3^2\cdot3=36\) 种，且自动满足 \(1\)，\(2\) 号中至少有 \(1\) 名新队员. 当有 \(2\) 名老队员和 \(1\) 名新队员时，选新队员有 \(3\) 种，新队员不能放在 \(3\) 号位，故有 \(2\) 个位置可放，两个老队员排列有 \(2\) 种，共有 \(3\cdot2\cdot2=12\) 种. 因此共有 \(36+12=48\) 种.
\end{solution}

\begin{problem}
若一条直线与一个正四棱柱各个面所成的角都为 \(\alpha\)，则 \(\cos \alpha = \)\fillinblank{}.
\end{problem}

\begin{answer}
\(\frac{\sqrt{6}}{3}\)
\end{answer}

\begin{solution}
设该直线的方向向量为 \(\bs{d}=(u,v,w)\). 正四棱柱的三个方向的侧面法向量可以分别取为 \((1,0,0)\)，\((0,1,0)\)，\((0,0,1)\). 直线与平面的夹角为 \(\alpha\) 时，方向向量与该平面法向量夹角的正弦为 \(\sin\alpha\)，所以 \(\frac{|u|}{|\bs{d}|}=\frac{|v|}{|\bs{d}|}=\frac{|w|}{|\bs{d}|}=\sin\alpha\). 由 \(u^2+v^2+w^2=|\bs{d}|^2\)，得 \(3\sin^2\alpha=1\). 由于夹角取锐角，\(\cos\alpha=\sqrt{1-\sin^2\alpha}=\sqrt{\frac23}=\frac{\sqrt6}{3}\).
\end{solution}

\section{解答题}
\begin{problem}
已知函数 \(f(x) = \sin^2 x + 2\sin x\cos x + 3 \cos^2 x\)，\(x \in \R\). 求：

\begin{enumerate}
\item 函数 \( f(x) \) 的最大值及取得最大值的自变量 \( x \) 的集合；
\item 函数 \( f(x) \) 的单调增区间.
\end{enumerate}
\end{problem}

\begin{solution}
\begin{enumerate}
\item 利用二倍角公式，得
\[
f(x)=2+\cos 2x+\sin 2x=2+\sqrt{2}\sin\left(2x+\frac{\pi}{4}\right)
\]
由于正弦函数的值不超过 \(1\)，所以 \(f(x)\) 的最大值为 \(2+\sqrt{2}\). 取得最大值时
\[
2x+\frac{\pi}{4}=\frac{\pi}{2}+2k\pi
\]
即 \(x=\frac{\pi}{8}+k\pi\)，其中 \(k\in\Z\).

\item \(f'(x)=2\cos 2x-2\sin 2x=2\sqrt{2}\cos\left(2x+\frac{\pi}{4}\right)\). 因而
\[
f'(x)>0\Longleftrightarrow -\frac{\pi}{2}+2k\pi<2x+\frac{\pi}{4}<\frac{\pi}{2}+2k\pi
\]
即 \(-\frac{3\pi}{8}+k\pi<x<\frac{\pi}{8}+k\pi\)，其中 \(k\in\Z\). 所以函数的单调增区间为
\[
\left(-\frac{3\pi}{8}+k\pi,\frac{\pi}{8}+k\pi\right)
\]
其中 \(k\in\Z\).
\end{enumerate}
\end{solution}

\begin{problem}
已知正方形 \(ABCD\)，\(E\)，\(F\) 分别是边 \(AB\)，\(CD\) 的中点，将 \(\triangle ADE\) 沿 \(DE\) 折起，如图所示，记二面角 \(A - DE - C\) 的大小为 \(\theta\)（\(0 < \theta < \pi\)）.

\begin{enumerate}
\item 证明：\(BF \parallel\) 平面 \(ADE\)；
\item 若 \(\triangle ACD\) 为正三角形，试判断点 \(A\) 在平面 \(BCDE\) 内的射影 \(G\) 是否在直线 \(EF\) 上. 证明你的结论，并求角 \(\theta\) 的余弦值.
\end{enumerate}

\FigureLayoutDeclare{img/2006/liaoning_science/q18_fig1.png}{13.0cm}{10bp 49bp 27bp 26bp}
\bitmapfigure[max width=0.82\linewidth]{img/2006/liaoning_science/q18_fig1.png}
\end{problem}

\begin{solution}
\begin{enumerate}
\item 设正方形边长为 \(a\)，在平面 \(BCDE\) 中建立坐标系，使 \(D=(0,0,0)\)，\(C=(a,0,0)\)，\(B=(a,a,0)\)，\(E=\left(\frac{a}{2},a,0\right)\)，\(F=\left(\frac{a}{2},0,0\right)\). 于是 \(\overrightarrow{BF}=\left(-\frac{a}{2},-a,0\right)=-\overrightarrow{DE}\). 直线 \(BF\) 与直线 \(DE\) 平行；又因 \(0<\theta<\pi\)，平面 \(ADE\) 与平面 \(BCDE\) 的交线为 \(DE\)，而 \(BF\subset BCDE\) 且 \(BF\ne DE\)，故 \(BF\not\subset ADE\)，从而 \(BF\parallel\) 平面 \(ADE\).

\item 设折起后 \(A=(u,v,w)\). 因为 \(\triangle ACD\) 为正三角形，且 \(AD=CD=a\)，又由折叠前后长度不变知 \(AE=\frac{a}{2}\)，所以
\[
\begin{cases}
u^2+v^2+w^2=a^2,\\
(u-a)^2+v^2+w^2=a^2,\\
\left(u-\frac{a}{2}\right)^2+(v-a)^2+w^2=\frac{a^2}{4}
\end{cases}
\]
前两个等式相减得 \(u=\frac{a}{2}\). 将其代入后两个等式，得 \(v=\frac{3a}{4}\)，\(w^2=\frac{3a^2}{16}\).
因此点 \(A\) 在平面 \(BCDE\) 上的射影为 \(G=\left(\frac{a}{2},\frac{3a}{4},0\right)\). 直线 \(EF\) 的方程为 \(x=\frac{a}{2}\)，且 \(0\leqslant y\leqslant a\)，所以 \(G\) 在线段 \(EF\) 上.

取 \(\bs{d}=\overrightarrow{DE}=\left(\frac{a}{2},a,0\right)\)，\(\bs{r}=\overrightarrow{DA}=\left(\frac{a}{2},\frac{3a}{4},w\right)\)，\(\bs{c}=\overrightarrow{DC}=(a,0,0)\).
将 \(\bs{r}\)，\(\bs{c}\) 分别在垂直于 \(DE\) 的方向上作正射影，得到
\(\bs{r}_{\perp}=\bs{r}-\frac{\bs{r}\cdot\bs{d}}{|\bs{d}|^2}\bs{d}=\left(\frac{a}{10},-\frac{a}{20},w\right)\)，\(\bs{c}_{\perp}=\bs{c}-\frac{\bs{c}\cdot\bs{d}}{|\bs{d}|^2}\bs{d}=\left(\frac{4a}{5},-\frac{2a}{5},0\right)\).
由 \(w^2=\frac{3a^2}{16}\)，有
\(\bs{r}_{\perp}\cdot\bs{c}_{\perp}=\frac{a^2}{10}\)，\(|\bs{r}_{\perp}|^2=\frac{a^2}{5}\)，\(|\bs{c}_{\perp}|^2=\frac{4a^2}{5}\).
故二面角 \(A-DE-C\) 的平面角 \(\theta\) 满足
\[
\cos\theta=\frac{\bs{r}_{\perp}\cdot\bs{c}_{\perp}}{|\bs{r}_{\perp}|\,|\bs{c}_{\perp}|}=\frac{1}{4}
\]
\end{enumerate}
\end{solution}

\begin{problem}
现有甲，乙两个项目，对甲项目每投资十万元，一年后利润是 \(1.2\) 万元，\(1.18\) 万元，\(1.17\) 万元的概率分别为 \(\frac{1}{6}\)，\(\frac{1}{2}\)，\(\frac{1}{3}\)；已知乙项目的利润与产品价格的调整有关. 在每次调整中，价格下降的概率都是 \(p\)（\(0 < p < 1\)），设乙项目产品价格在一年内进行 \(2\) 次独立的调整，记乙项目产品价格在一年内的下降次数为 \(\xi\)，对乙项目每投资十万元，\(\xi\) 取 \(0\)，\(1\)，\(2\) 时，一年后相应利润是 \(1.3\) 万元，\(1.25\) 万元，\(0.2\) 万元. 随机变量 \(\xi_1\)，\(\xi_2\) 分别表示对甲，乙两项目各投资十万元一年的利润.

\begin{enumerate}
\item 求 \( \xi_1 \)，\( \xi_2 \) 的概率分布和数学期望 \( E(\xi_1) \)，\( E(\xi_2) \)；
\item 当 \( E(\xi_1) < E(\xi_2) \) 时，求 \(p\) 的取值范围.
\end{enumerate}
\end{problem}

\begin{solution}
\begin{enumerate}
\item 随机变量 \(\xi_1\) 的分布为
\(P(\xi_1=1.2)=\frac{1}{6}\)，\(P(\xi_1=1.18)=\frac{1}{2}\)，\(P(\xi_1=1.17)=\frac{1}{3}\).
两次价格调整相互独立，故 \(\xi\) 服从参数为 \(2\)，\(p\) 的二项分布，从而
\(P(\xi=0)=(1-p)^2\)，\(P(\xi=1)=2p(1-p)\)，\(P(\xi=2)=p^2\).
根据利润与 \(\xi\) 的对应关系，随机变量 \(\xi_2\) 的分布为
\(P(\xi_2=1.3)=(1-p)^2\)，\(P(\xi_2=1.25)=2p(1-p)\)，\(P(\xi_2=0.2)=p^2\).
所以
\[
E(\xi_1)=1.2\times\frac{1}{6}+1.18\times\frac{1}{2}+1.17\times\frac{1}{3}=1.18
\]
\[
E(\xi_2)=1.3(1-p)^2+1.25\times2p(1-p)+0.2p^2=1.3-0.1p-p^2
\]

\item 由 \(E(\xi_1)<E(\xi_2)\)，得
\[
1.18<1.3-0.1p-p^2
\Longleftrightarrow p^2+0.1p-0.12<0
\Longleftrightarrow -0.4<p<0.3
\]
再结合 \(0<p<1\)，可得 \(p\) 的取值范围为 \(0<p<\frac{3}{10}\).
\end{enumerate}
\end{solution}

\begin{problem}
已知点 \(A(x_{1},y_{1})\)，\(B(x_{2},y_{2})\)（\(x_{1}\)，\(x_{2}\neq 0\)）是抛物线 \(y^{2} = 2px\)（\(p > 0\)）上的两个动点，\(O\) 是坐标原点，向量 \(\overrightarrow{OA}\)，\(\overrightarrow{OB}\) 满足 \(\left|\overrightarrow{OA} +\overrightarrow{OB}\right| = \left|\overrightarrow{OA} -\overrightarrow{OB}\right|\). 设圆 \(C\) 的方程为 \(x^{2} + y^{2} - (x_{1} + x_{2})x - (y_{1} + y_{2})y = 0\).

\begin{enumerate}
\item 证明线段 \(AB\) 是圆 \(C\) 的直径；
\item 当圆 \(C\) 的圆心到直线 \(x - 2y = 0\) 的距离的最小值为 \(\frac{2\sqrt{5}}{5}\) 时，求 \(p\) 的值.
\end{enumerate}
\end{problem}

\begin{solution}
\begin{enumerate}
\item 由已知条件，平方后相减可得
\[
\left|\overrightarrow{OA}+\overrightarrow{OB}\right|^2-\left|\overrightarrow{OA}-\overrightarrow{OB}\right|^2=4\overrightarrow{OA}\cdot\overrightarrow{OB}=0
\]
因此 \(x_1x_2+y_1y_2=0\). 将圆 \(C\) 的方程配方，得
\[
\left(x-\frac{x_1+x_2}{2}\right)^2+\left(y-\frac{y_1+y_2}{2}\right)^2=\frac{(x_1+x_2)^2+(y_1+y_2)^2}{4}
\]
所以圆心为线段 \(AB\) 的中点. 又因为
\[
\left|\overrightarrow{OA}+\overrightarrow{OB}\right|=\left|\overrightarrow{OA}-\overrightarrow{OB}\right|=|\overrightarrow{AB}|
\]
圆 \(C\) 的半径等于 \(\frac{1}{2}|\overrightarrow{AB}|\). 故线段 \(AB\) 是圆 \(C\) 的直径.

\item 设抛物线上点的参数为 \(t\)，则点的坐标可写为 \(\left(\frac{p}{2}t^2,pt\right)\). 令 \(A\)，\(B\) 对应的参数分别为 \(t_1\)，\(t_2\). 因为 \(x_1\)，\(x_2\neq0\)，所以 \(t_1\)，\(t_2\neq0\). 由 \(\overrightarrow{OA}\perp\overrightarrow{OB}\)，得
\[
\frac{p^2}{4}t_1^2t_2^2+p^2t_1t_2=0
\]
即 \(t_1t_2=-4\). 令 \(t_1=t\)，则 \(t_2=-\frac{4}{t}\).

圆心 \(M\) 到直线 \(x-2y=0\) 的距离为
\[
d=\frac{|x_1+x_2-2y_1-2y_2|}{2\sqrt{5}}
\]
代入参数表达式，得
\[
d=\frac{p}{4\sqrt{5}}\left|t^2+\frac{16}{t^2}-4t+\frac{16}{t}\right|
\]
令 \(u=t-\frac{4}{t}\)，则
\[
t^2+\frac{16}{t^2}-4t+\frac{16}{t}=u^2+8-4u=(u-2)^2+4\geqslant4
\]
当 \(t-\frac{4}{t}=2\) 时等号成立，此方程有非零实根 \(t=1\pm\sqrt{5}\)，故最小距离为 \(\frac{p}{\sqrt{5}}\). 由题设
\[
\frac{p}{\sqrt{5}}=\frac{2\sqrt{5}}{5}=\frac{2}{\sqrt{5}}
\]
所以 \(p=2\).
\end{enumerate}
\end{solution}

\begin{problem}
已知函数 \( f(x) = \frac{1}{3} ax^3 + bx^2 + cx + d \)，其中 \(a\)，\(b\)，\(c\) 是以 \(d\) 为公差的等差数列，且 \(a > 0\)，\(d > 0\). 设 \(x_0\) 为 \(f(x)\) 的极小值点. 在 \(\left[1 - \frac{2b}{a}, 0\right]\) 上，\(f'(x)\) 在 \(x_1\) 处取得最大值，在 \(x_2\) 处取得最小值. 将 \((x_0, f(x_0))\)，\((x_1, f'(x_1))\)，\((x_2, f'(x_2))\) 依次记为 \(A\)，\(B\)，\(C\).

\begin{enumerate}
\item 求 \(x_0\)；
\item 若 \(\triangle ABC\) 有一边平行于 \(x\) 轴，且面积为 \(2 + \sqrt{3}\)，求 \(a\)，\(d\) 的值.
\end{enumerate}
\end{problem}

\begin{solution}
\begin{enumerate}
\item 因为 \(a\)，\(b\)，\(c\) 是以 \(d\) 为公差的等差数列，由等差数列定义，\(b=a+d\)，\(c=a+2d\).
于是
\[
f'(x)=ax^2+2bx+c=ax^2+2(a+d)x+a+2d=(x+1)\left(ax+a+2d\right)
\]
其两个零点为 \(-1-\frac{2d}{a}\) 与 \(-1\). 因为 \(a>0\)，所以 \(f'(x)\) 在 \(x=-1\) 处由负变正，故 \(f(x)\) 的极小值点为 \(x_0=-1\).

\item 区间 \(\left[1-\frac{2b}{a},0\right]=\left[-1-\frac{2d}{a},0\right]\). 函数 \(f'(x)\) 是开口向上的二次函数，其顶点横坐标为
\[
x_2=-\frac{b}{a}=-1-\frac{d}{a}
\]
又 \(f'\left(-1-\frac{2d}{a}\right)=0\)，而 \(f'(0)=c=a+2d>0\)，所以最大值在右端点取得，即 \(x_1=0\)，最小值在顶点取得，且
\[
f'(x_2)=-\frac{d^2}{a}
\]
从而
\(A=\left(-1,-\frac{a}{3}\right)\)，\(B=(0,a+2d)\)，\(C=\left(-1-\frac{d}{a},-\frac{d^2}{a}\right)\).
其中 \(y_B=a+2d>0\)，而 \(y_A=-\frac{a}{3}<0\)，\(y_C=-\frac{d^2}{a}<0\). 因此若三角形 \(ABC\) 有一边平行于 \(x\) 轴，只能是边 \(AC\)，于是
\(-\frac{a}{3}=-\frac{d^2}{a}\)，\(d=\frac{a}{\sqrt{3}}\).
此时 \(|AC|=\frac{d}{a}=\frac{1}{\sqrt{3}}\)，点 \(B\) 到直线 \(AC\) 的距离为
\(a+2d+\frac{a}{3}=\frac{4a}{3}+2d\).
由三角形面积为 \(2+\sqrt{3}\)，得
\[
\frac{1}{2}\cdot\frac{1}{\sqrt{3}}\left(\frac{4a}{3}+\frac{2a}{\sqrt{3}}\right)=2+\sqrt{3}
\]
即 \(\frac{a(2+\sqrt{3})}{3\sqrt{3}}=2+\sqrt{3}\). 所以 \(a=3\sqrt{3}\)，从而 \(d=3\).
\end{enumerate}
\end{solution}

\begin{problem}
已知 \(f_{0}(x) = x^{n}\)，\(f_{k}(x) = \frac{f'_{k-1}(x)}{f_{k-1}(1)}\)，其中 \(k \leqslant n\)（\(n\)，\(k\in\mathbf{N}_{+}\)）. 设 \(F(x) = \mathrm C_n^0 f_0(x^{2}) + \mathrm C_n^1 f_1(x^{2}) + \cdots + \mathrm C_n^n f_n(x^{2})\)，\(x \in [-1,1]\).

\begin{enumerate}
\item 写出 \(f_k(1)\)；
\item 证明：对任意的 \(x_{1}\)，\(x_{2}\in[-1,1]\)，恒有 \(|F(x_{1}) - F(x_{2})| \leqslant 2^{n-1}(n+2) - n - 1\).
\end{enumerate}
\end{problem}

\begin{solution}
\begin{enumerate}
\item 由 \(f_0(x)=x^n\) 得 \(f_0(1)=1\). 先求得
\[
f_1(x)=nx^{n-1}
\]
若 \(2\leqslant k\leqslant n\)，并且 \(f_{k-1}(x)=(n-k+2)x^{n-k+1}\)，则
\[
f_k(x)=\frac{f'_{k-1}(x)}{f_{k-1}(1)}=(n-k+1)x^{n-k}
\]
由数学归纳法，\(f_k(x)=(n-k+1)x^{n-k}\)，所以 \(f_0(1)=1\)，\(f_k(1)=n-k+1\)（\(1\leqslant k\leqslant n\)）.

\item 令 \(u=x^2\)，则 \(0\leqslant u\leqslant1\)，且
\[
F(x)=P(u)=u^n+\sum_{k=1}^{n}\mathrm C_n^k(n-k+1)u^{n-k}
\]
这个多项式的各项系数均为非负数，因此当 \(0\leqslant u_1\leqslant u_2\leqslant1\) 时，有 \(P(u_1)\leqslant P(u_2)\). 对任意 \(x_1\)，\(x_2\in[-1,1]\)，取 \(u_i=x_i^2\)，便有
\[
|F(x_1)-F(x_2)|=|P(u_1)-P(u_2)|\leqslant P(1)-P(0)
\]
再由二项式系数的和及加权和
\[
\sum_{k=0}^{n}\mathrm C_n^k(n-k+1)=n2^{n-1}+2^n=2^{n-1}(n+2)
\]
注意到上式中 \(k=0\) 项为 \(n+1\)，而 \(P(1)\) 的对应项为 \(1\)，故
\[
P(1)=2^{n-1}(n+2)-n
\]
当 \(u=0\) 时只有 \(k=n\) 项不为零，所以 \(P(0)=1\). 因而
\[
|F(x_1)-F(x_2)|\leqslant P(1)-P(0)=2^{n-1}(n+2)-n-1
\]
这就证明了所要的不等式.
\end{enumerate}
\end{solution}
